\renewcommand{\currentchapter}{Referensi}

\begingroup
\setbeamercolor{background canvas}{bg=StudentNavy}
\begin{frame}[plain]
  \vspace{0.65cm}
  \kicker{Laporan UTBK 2027}
  \vspace{0.08cm}
  {\color{StudentWhite}\bfseries\fontsize{28}{33}\selectfont Daftar Pustaka\par}
  \vspace{0.38cm}
  {\color{StudentSoft}\fontsize{14.5}{18}\selectfont Seluruh sumber resmi, publik, dan internal yang digunakan dalam dokumen}\par
  \vfill
  {\color{StudentGold}\rule{2.3cm}{2.5pt}}\par
  \vspace{0.20cm}
  {\color{StudentWhite}\bfseries\fontsize{15}{18}\selectfont Balya Rochmadi}\par
  {\color{StudentSoft}\fontsize{11.5}{14}\selectfont SMAN 1 Muntilan bersama StudentPro}
  \vspace{0.42cm}
\end{frame}
\endgroup

\begin{frame}{Catatan status sumber}
  \kicker{Cara Membaca Bibliografi}
  \begin{columns}[T,onlytextwidth]
    \begin{column}{0.30\textwidth}
      \begin{block}{Resmi}
        Dokumen dan laman Panitia SNPMB menjadi acuan struktur UTBK 2026.
      \end{block}
    \end{column}
    \begin{column}{0.30\textwidth}
      \begin{block}{Publik}
        Pemberitaan digunakan sebagai sumber sekunder dan tidak menggantikan dokumen resmi.
      \end{block}
    \end{column}
    \begin{column}{0.30\textwidth}
      \begin{block}{Internal}
        LPJ, rekap, bank materi, serta catatan StudentPro digunakan untuk evaluasi dan perencanaan program.
      \end{block}
    \end{column}
  \end{columns}
  \vspace{0.26cm}
  \begin{beamercolorbox}[rounded=true,sep=0.18cm,wd=\textwidth]{block body}
    \centering\bfseries Catatan internal yang belum lengkap tidak diperlakukan sebagai ketentuan resmi atau klaim publik.
  \end{beamercolorbox}
\end{frame}

\nocite{*}
\begin{frame}[allowframebreaks]{Daftar pustaka}
  \kicker{Rujukan Terpusat}
  {\fontsize{9.2}{11.5}\selectfont
  \bibliographystyle{apalike}
  \bibliography{referensi}}
\end{frame}
